% Abstract — 100-250 words (Handbook 8.1.1)
%
% === Questions for Martin (meeting notes) ===
%
% 1. Lead with privacy or engineering? The project is both a privacy system
%    and a hard engineering challenge (building on immature Veilid). Which
%    angle do markers care about more?
%
% 2. Seller-only-gets-a-route: the seller never learns the winner's identity
%    — just an ephemeral onion route. Stronger than most sealed-bid systems.
%    Worth emphasising in the abstract or too in-the-weeds?
%
% 3. Monero descope framing. Time went to MPC tunnel hardening + route
%    recovery + E2E tests instead. Was this the right tradeoff?
%
% 4. LSEPI / dual-use honesty level. The system *could* be deployed as a
%    real anonymous marketplace. The limitations (no Sybil resistance, seller
%    trust, liveness) are all fixable engineering problems, not fundamental.
%    How candid should the report be about this?
%
% 5. Assumptions up front vs in Results. Ch1 currently plans to state what
%    IS private, what ISN'T, trust assumptions, liveness constraints. Too
%    much for an intro, or good that markers see it immediately?
%
% 6. OT descope. PDD said 1-out-of-N OT for decryption key delivery.
%    Implemented commitment+challenge instead. Seller already sees winning
%    bid via MPC (party 0), so no additional info leak. Discuss as limitation
%    or pragmatic equivalent?
%
% ===========================================
\section*{Abstract}
\addcontentsline{toc}{section}{Abstract}

Bidder privacy is absent from mainstream online auctions.  Open-bid platforms
reveal the bidding process to all participants; sealed-bid schemes hide bid
values but still require a trusted auctioneer who sees every bid.  In both
models, losing bidders' information is unnecessarily exposed, and metadata ---
who bid, when, how often --- is visible to platform operators and network
observers.  This project explored whether a fully decentralised auction could
protect bidder privacy without any trusted intermediary.

A peer-to-peer sealed-bid marketplace was built on Veilid, a privacy-first
application framework providing onion-routed messaging, a secure distributed
hash table, and public-key identities with no central server.  Bids are
cryptographically committed and the winner is determined jointly by all
participants using the MASCOT multi-party computation protocol, so that losing
bid values are never revealed to anyone.  Network-level privacy is provided by
Veilid's private routing, which obscures participants' IP addresses and
locations.  Even the seller, who learns the winning bid value, receives only an
ephemeral onion route to the winner --- not their identity.

The prototype was implemented in Rust (approximately 14,800 lines of original
code) and tested on a 20-node devnet with automated end-to-end auction
scenarios.  A security analysis identified residual privacy limitations: bid
timing can leak party assignments, participation metadata is visible on the
DHT, and any single party going offline during MPC execution prevents the
auction from completing.


% Scope the definitions to be strict!
% Metadata about bids / What are you trying to protect, who are you trying to protect it from?`